\chapter{Conjuntos Finitos, Enumerables y No Enumerables}

Una vez que tenemos los conceptos generales y los hechos básicos sobre conjuntos y funciones, los aplicaremos, en este capítulo, para estudiar las nociones de conjunto finito y conjunto infinito.

A Cantor se debe el descubrimiento fundamental de que hay diversos tipos de infinito, así como el análisis de esos tipos.

La noción de conjunto enumerable está estrictamente ligada al conjunto \(\mathbb{N}\) de los números naturales. Por ello iniciaremos este capítulo con una breve presentación de la teoría de los números naturales a partir de los axiomas de Peano.

Los axiomas de Peano exhiben los números naturales como "números ordinales", es decir, objetos que ocupan lugares determinados en una secuencia ordenada: 1 es el primer número natural, 2 es el número que viene justo después del 1, 3 viene a continuación del 2, etc. Pero los números naturales también pueden ocurrir como “Números cardinales”, es decir, como resultado de una operación de conteo, en respuesta a la pregunta: ¿cuántos elementos tiene este conjunto?

En el §2, empleamos los números naturales para el conteo de los conjuntos finitos, mostrando cómo pueden ser considerados como números cardinales y completando, por tanto, su descripción.

Una exposición sistemática de los sistemas numéricos utilizados en el Análisis Matemático puede hacerse a partir de los números naturales, mediante sucesivas extensiones del concepto de número: primero se amplía \(\mathbb{N}\) para obtener el conjunto \(\mathbb{Z}\) de los números enteros; luego se extiende \(\mathbb{Z}\), pasando al conjunto \(\mathbb{Q}\) de los números racionales; de este se pasa al conjunto \(\mathbb{R}\) de los reales y, de ahí, al conjunto \(\mathbb{C}\) de los complejos. Esta elaboración, aunque instructiva, es un proceso lento.

\section{Números naturales}

Toda la teoría de los números naturales puede deducirse de los tres axiomas siguientes, conocidos como axiomas de Peano.

Se dan, como objetos no definidos, un conjunto \(\mathbb{N}\), cuyos elementos se llaman números naturales, y una función \(s: \mathbb{N} \to \mathbb{N}\). Para cada \(n \in \mathbb{N}\), el número \(s(n)\), valor que la función asume en el punto \(n\), se llama el sucesor de \(n\).

La función \(s\) satisface los siguientes axiomas:

P1. \(s: \mathbb{N} \to \mathbb{N}\) es inyectiva. En otros términos: \(m, n \in \mathbb{N}\), \(s(m) = s(n) \Rightarrow m = n\). O, en palabras, dos números que tienen el mismo sucesor son iguales.

P2. \(\mathbb{N} - s(\mathbb{N})\) consta de un solo elemento. Es decir, existe un único número natural que no es sucesor de ningún otro. Ese se llama “uno” y se representa con el símbolo 1. Así, cualquiera que sea \(n \in \mathbb{N}\), se tiene \(1 \neq s(n)\). Por otro lado, si \(n \neq 1\) entonces existe un (único) \(n_0 \in \mathbb{N}\) tal que \(s(n_0) = n\).

P3. (Principio de Inducción). Si \(X \subset \mathbb{N}\) es un subconjunto tal que \(1 \in X\) y, para todo \(n \in X\) se tiene también \(s(n) \in X\), entonces \(X = \mathbb{N}\).

El Principio de Inducción puede también enunciarse de la siguiente manera:

Sea \(\mathcal{P}\) una propiedad referente a números naturales. Si 1 goza de la propiedad \(\mathcal{P}\) y si, del hecho de que un número natural \(n\) goce de \(\mathcal{P}\) se puede concluir que \(n + 1\) goza de la propiedad \(\mathcal{P}\), entonces todos los números naturales gozan de esa propiedad.

Una demostración en la que se emplea el axioma P3 se llama una demostración por inducción.

Para dar un ejemplo de demostración por inducción, mostremos que para todo \(n \in \mathbb{N}\) se tiene \(s(n) \neq n\). En efecto, sea \(X = \{n \in \mathbb{N}; s(n) \neq n\}\). Se tiene \(1 \in X\), pues 1 no es sucesor de ningún número, en particular \(1 \neq s(1)\). Además, \(n \in X \Rightarrow n \neq s(n) \Rightarrow\) (por la inyectividad de \(s\)) \(s(n) \neq s[s(n)] \Rightarrow s(n) \in X\). Así, \(n \in X \Rightarrow s(n) \in X\). Como \(1 \in X\), se sigue del axioma P3 que \(X = \mathbb{N}\), o sea \(n \neq s(n)\) para todo \(n \in \mathbb{N}\).

El Principio de Inducción es muy útil para demostrar proposiciones que se refieren a enteros. Él está implícito en todos los argumentos donde se dice “y así sucesivamente”, “y así seguidamente” o “etc.”.

No menos importante que demostrar proposiciones por inducción es saber definir objetos inductivamente.

Las definiciones por inducción se basan en la posibilidad de iterar una función \(f: X \to X\) un número arbitrario \(n\) de veces.

Más precisamente, sea \(f: X \to X\) una función cuyo dominio y codominio son el mismo conjunto \(X\). A cada \(n \in \mathbb{N}\) podemos, de modo único, asociar una función \(f^n: X \to X\) de tal manera que \(f^1 = f\) y \(f^{s(n)} = f \circ f^n\). En particular, si llamamos \(2 = s(1)\), \(3 = s(2)\), tendremos \(f^2 = f \circ f\), \(f^3 = f \circ f \circ f\).

En una exposición sistemática de la teoría de los números naturales, la existencia de la \(n\)-ésima iterada \(f^n\) de una función \(f: X \to X\) es un teorema, llamado “Teorema de la Definición por Inducción”. No daremos su demostración aquí. Solo observaremos que no nos sería posible, a estas alturas, definir \(f^n\) simplemente como \(f \circ f \circ \cdots \circ f\) (\(n\) veces) pues “\(n\) veces” es una expresión sin sentido en el contexto de los axiomas de Peano, ya que un número natural \(n\) es, por ahora, solo un elemento del conjunto \(\mathbb{N}\) (un número ordinal), sin condiciones de servir de respuesta a la pregunta “¿cuántas veces?”, hasta que se le otorgue la condición de número cardinal.

Admitamos, por tanto, que, dada una función \(f: X \to X\), sabemos asociar, de modo único, a cada número natural \(n \in \mathbb{N}\), una función \(f^n: X \to X\), llamada la \(n\)-ésima iterada de \(f\), de tal modo que \(f^1 = f\) y \(f^{s(n)} = f \circ f^n\).

Veamos un ejemplo de definición por inducción. Usando las iteradas de la función \(s: \mathbb{N} \to \mathbb{N}\), definiremos la suma de números naturales. Dados \(m, n \in \mathbb{N}\), su suma \(m + n \in \mathbb{N}\) se define por:

\[m + n = s^n(m).\]

Así, sumar \(m\) con 1 es tomar el sucesor de \(m\) mientras que, en general, sumar \(m\) con \(n\) es partir de \(m\) e iterar \(n\) veces la operación de tomar el sucesor. En otras palabras, tenemos, por definición:

\[m + 1 = s(m),\]

\[m + s(n) = s(m + n).\]

Así, si queremos, podremos prescindir de la notación \(s(n)\) para indicar el sucesor de \(n\) y usar la notación definitiva \(n + 1\) para representar dicho sucesor. Esto se hará gradualmente. Con la notación definitiva, la última de las igualdades anteriores se lee:

\[m + (n + 1) = (m + n) + 1.\]

Vemos que la propia definición de la suma \(m + n\) ya contiene una indicación de que goza de la propiedad asociativa. Probaremos ahora, en general, que se tiene:

\[m + (n + p) = (m + n) + p,\]

cualesquiera que sean \(m\), \(n\), \(p \in \mathbb{N}\).

La demostración de la propiedad asociativa se hace así: sea \(X\) el conjunto de los números naturales \(p\) tales que \(m + (n + p) = (m + n) + p\), cualesquiera que sean \(m\), \(n \in \mathbb{N}\). Ya vimos que \(1 \in X\). Además, si \(p \in X\), entonces,
\[m + (n + s(p)) = m + s(n + p) = s(m + (n + p)) = s((m + n) + p) = (m + n) + s(p).\]

(En el tercer signo de igualdad, usamos la hipótesis \(p \in X\) y en los demás la definición de suma.) Luego, \(p \in X \Rightarrow s(p) \in X\). Como \(1 \in X\), concluimos, por inducción, que \(X = \mathbb{N}\), es decir, \(m + (n + p) = (m + n) + p\), cualesquiera que sean \(m\), \(n\), \(p \in \mathbb{N}\).

Las propiedades formales de la adición se relacionan a continuación:

Asociatividad: \(m + (n + p) = (m + n) + p\);\\
Conmutatividad: \(m + n = n + m\);\\
Ley de cancelación: \(m + n = m + p \Rightarrow n = p\);\\
Tricotomía: dados \(m\), \(n \in \mathbb{N}\), exactamente una de las tres alternativas siguientes puede ocurrir: o \(m = n\), o existe \(p \in \mathbb{N}\) tal que \(m = n + p\), o bien existe \(q \in \mathbb{N}\) con \(n = m + q\).

Omitimos las demostraciones de esas propiedades, que se hacen por inducción.

La relación de orden entre los números naturales se define en términos de la adición. Dados los números naturales \(m\), \(n\) decimos que \(m\) es menor que \(n\) y escribimos
\[m < n,\]

para significar que existe \(p \in \mathbb{N}\) tal que \(n = m + p\). En las mismas condiciones, decimos que \(n\) es mayor que \(m\) y escribimos \(n > m\). La notación \(m \leq n\) significa que \(m\) es menor o igual que \(n\).

La relación \(<\) goza de las siguientes propiedades:

Transitividad: si \(m < n\) y \(n < p\) entonces \(m < p\).\\
Tricotomía: dados \(m\), \(n\), exactamente una de las siguientes alternativas puede ocurrir: o \(m = n\), o \(m < n\) o \(n < m\).\\
Monotonicidad de la adición: si \(m < n\) entonces, para todo \(p \in \mathbb{N}\) se tiene \(m + p < n + p\).

Probemos la última: \(m < n\) significa que existe \(q \in \mathbb{N}\) tal que \(m + q = n\). Entonces \((m + p) + q = n + p\), y, por tanto, \(m + p < n + p\).

Observemos que \(m < n\) significa que, para cierto \(p \in \mathbb{N}\), tenemos \(n = s^p(m)\), esto es, \(n\) es el sucesor del sucesor ... del sucesor de \(m\).

Introduciremos ahora la multiplicación de números naturales.

Así, como la suma \(m + n\) fue definida como el resultado que se obtiene cuando se itera \(n\) veces, a partir de \(m\), la operación de tomar el sucesor, definiremos el producto \(m \cdot n\) como la suma de \(n\) sumandos iguales a \(m\), o mejor, el resultado que se obtiene cuando se añade a \(m\), \(n - 1\) veces, el mismo número \(m\).

En términos precisos, para cada \(m \in \mathbb{N}\), sea \(f_m: \mathbb{N} \to \mathbb{N}\) la función definida por \(f_m(p) = p + m\). (Es decir, \(f_m\) es la función “sumar \(m\)”). Usaremos esta función para definir la multiplicación de números naturales.

El producto de dos números naturales se define así, \(m \cdot 1 = m\) y \(m \cdot (n + 1) = (f_m)^n(m)\).

En palabras: multiplicar un número \(m\) por 1 no lo altera. Multiplicar \(m\) por un número mayor que 1, es decir, por un número de la forma \(n+1\), es iterar \(n\) veces la operación de sumar \(m\), comenzando con \(m\). Así, por ejemplo, \(m \cdot 2 = f_m(m) = m+m\), \(m \cdot 3 = (f_m)^2(m) = m + m + m\).

Recordando la definición de \((f_m)^n\), vemos que el producto \(m \cdot n\) está definido inductivamente por las propiedades siguientes:
\[
m \cdot 1 = m,
\]
\[
m \cdot (n + 1) = m \cdot n + m.
\]

La última igualdad anterior ya sugiere que el producto debe gozar de la propiedad distributiva
\[
m \cdot (n + p) = m \cdot n + m \cdot p.
\]

Demostremos este hecho. Sea \(X\) el conjunto de los números \(p \in \mathbb{N}\) tales que \((m + n) \cdot p = m \cdot p + n \cdot p\) cualesquiera que sean \(m, n \in \mathbb{N}\). Acabamos de observar que \(1 \in X\). Además, si \(p \in X\), concluiremos que

\[
(m + n) \cdot (p + 1) = (m + n) \cdot p + m + n
\]
\[
= m \cdot p + n \cdot p + m + n = m \cdot p + m + n \cdot p + n
\]
\[
= m \cdot (p + 1) + n \cdot (p + 1).
\]

(En las igualdades anteriores usamos, sucesivamente, la definición de producto, la hipótesis de que \(p \in X\), la asociatividad y la conmutatividad de la adición y, nuevamente, la definición de producto).

Se sigue que \(p + 1 \in X\). Así, \(X = \mathbb{N}\), es decir, \((m + n) \cdot p = m \cdot p + n \cdot p\), cualesquiera que sean \(m, n, p \in \mathbb{N}\).

Las principales propiedades de la multiplicación son:

Asociatividad: \(m \cdot (n \cdot p) = (m \cdot n) \cdot p\);\\
Conmutatividad: \(m \cdot n = n \cdot m\);\\
Ley de cancelación: \(m \cdot p = n \cdot p \Rightarrow m = n\);\\
Distributividad: \(m(n + p) = m \cdot n + m \cdot p\);\\
Monotonicidad: \(m < n \Rightarrow m \cdot p < n \cdot p\).

\section{Buena ordenación y el Segundo Principio de Inducción}

Sea \( X \) un conjunto de números naturales. Se dice que un número \( p \in X \) es el menor elemento de \( X \) (o elemento mínimo de \( X \)) cuando se tiene \( p \leq n \) para todo \( n \in X \). Por ejemplo, 1 es el menor elemento del conjunto \( \mathbb{N} \) de todos los números naturales. Con mayor razón, cualquiera que sea \( X \subset \mathbb{N} \) con \( 1 \in X \), 1 es el menor elemento de \( X \).

Dado \( X \subset \mathbb{N} \), si \( p \in X \) y \( q \in X \) son ambos los menores elementos de \( X \), entonces \( p \leq q \) y \( q \leq p \), de donde \( p = q \). Así, el menor elemento de un conjunto es único.

Análogamente, si \( X \subset \mathbb{N} \), un número \( p \in X \) se llama el mayor elemento de \( X \) (o elemento máximo de \( X \)) cuando se tiene \( p \geq n \) para todo \( n \in X \). No todo conjunto de números naturales posee un elemento máximo. Por ejemplo, el propio \( \mathbb{N} \) no tiene un mayor elemento ya que, para todo \( n \in \mathbb{N} \), se tiene \( n + 1 > n \).

Si existe el elemento máximo de un conjunto \( X \subset \mathbb{N} \), es único. En efecto, si \( p \in X \) y \( q \in X \) son ambos máximos, entonces \( p \geq q \) y \( q \geq p \), de donde \( p = q \).

Un resultado de gran importancia, incluso como método de demostración, es el hecho de que todo conjunto no vacío de números naturales posee un menor elemento. Este hecho se conoce como el Principio de Buena Ordenación.

\textbf{Teorema 1}. (Principio de Buena Ordenación). Todo subconjunto no vacío \( A \subset \mathbb{N} \) posee un elemento mínimo.

\textbf{Demostración}. Usando la notación \( I_n = \{p \in \mathbb{N} : 1 \leq p \leq n\} \), consideremos el conjunto \( X \subset \mathbb{N} \) formado por los números \( n \in \mathbb{N} \) tales que \( I_n \subset \mathbb{N} - A \). (Así, decir que \( n \in X \) significa afirmar que \( n \notin A \) y que todos los números naturales menores que \( n \) tampoco pertenecen a \( A \).) Si tenemos \( 1 \in A \), el teorema estará demostrado pues 1 será el menor elemento de \( A \). Si, en cambio, es \( 1 \notin A \), entonces \( 1 \in X \). Por otro lado, tenemos \( X \neq \mathbb{N} \) (pues \( X \subset \mathbb{N} - A \) y \( A \neq \emptyset \)). Así, \( X \) cumple la primera parte de la hipótesis de P3 (contiene 1) pero no satisface la conclusión de P3 (no es igual a \(\mathbb{N}\)). Luego no puede cumplir la segunda parte de la hipótesis. Esto quiere decir: debe existir algún \( n \in X \) tal que \( n + 1 \notin X \). Sea \( a = n + 1 \). Entonces todos los enteros desde 1 hasta \( n \) pertenecen al complemento de \( A \) pero \( a = n + 1 \) pertenece a \( A \). De esta manera, \( a \) es el menor elemento del conjunto \( A \), lo que demuestra el teorema.

Del Principio de Buena Ordenación se deduce una proposición conocida como el Segundo Principio de Inducción, que demostramos ahora.

\textbf{Teorema 2}. (Segundo Principio de Inducción.) Sea \( X \subset \mathbb{N} \) un conjunto con la siguiente propiedad: dado \( n \in \mathbb{N} \), si \( X \) contiene todos los números naturales \( m \) tales que \( m < n \), entonces \( n \in X \). En estas condiciones, \( X = \mathbb{N} \).

\textbf{Demostración}. Sea \( Y = \mathbb{N} - X \). Afirmamos que \( Y = \emptyset \). En efecto, si \( Y \) no fuera vacío, existiría un elemento mínimo \( p \in Y \). Entonces, para todo número natural \( m < p \), sería \( m \in X \). Por la hipótesis hecha sobre \( X \), tendríamos \( p \in X \), una contradicción.

El Segundo Principio de Inducción constituye un método útil para demostración de proposiciones referentes a números naturales. También puede enunciarse así: Sea \( \mathcal{P} \) una propiedad relativa a números naturales. Si, dado \( n \in \mathbb{N} \), del hecho de que todo número natural \( m < n \) goce de la propiedad \( \mathcal{P} \) puede inferirse que \( n \) goza de la propiedad \( \mathcal{P} \), entonces todo número natural goza de \( \mathcal{P} \).

Ejemplo 1. Un número natural \( p \) se llama primo cuando \( p \neq 1 \) y no se puede escribir \( p = m \cdot n \) con \( m < p \) y \( n < p \). El llamado Teorema Fundamental de la Aritmética dice que todo número natural se descompone, de modo único, como producto de factores primos. La demostración utiliza el Segundo Principio de Inducción. En efecto, dado \( n \in \mathbb{N} \), supongamos que todo número natural menor que \( n \) puede descomponerse como producto de factores primos. Entonces, o \( n \) es primo (y en este caso \( n \) es, de modo trivial, un producto de factores primos) o bien \(n = m \cdot k\), con \(m < n\) y \(k < n\). Por la hipótesis de inducción, \(m\) y \(k\) son productos de factores primos. Se sigue que \(n\) también lo es. Por el Segundo Principio de Inducción, concluimos que todo número natural es producto de números primos.

Concluiremos este parágrafo hablando del método general de definición por inducción.

Sea \(X\) un conjunto cualquiera. Se trata de definir una función \(f: \mathbb{N} \to X\). Supongamos que se nos da el valor \(f(1)\) y se da también una regla que nos permita obtener \(f(n)\) siempre que conozcamos los valores \(f(m)\), para todo \(m < n\). Entonces existe una, y solo una, función \(f: \mathbb{N} \to X\) en esas condiciones. Este es el método general de definición por inducción (o por recurrencia).

Ejemplo 2. Fijado \(\alpha \in \mathbb{N}\), definamos una función \(f: \mathbb{N} \to \mathbb{N}\), inductivamente, poniendo \(f(1) = \alpha\) y \(f(n+1) = \alpha \cdot f(n)\). La función \(f\) cumple entonces \(f(2) = \alpha \cdot \alpha\), \(f(3) = \alpha \cdot \alpha \cdot \alpha\), etc. Luego \(f(n) = \alpha^n\). Acabamos de definir, por inducción, la \(n\)-ésima potencia del número natural \(\alpha\). Nótese que, en este caso, tenemos solo la iteración de la multiplicación por \(\alpha\), es decir, la forma más simple de definición por inducción. Los ejemplos siguientes requieren la forma general, esto es, no se reducen inmediatamente a la iteración de una función dada.

Ejemplo 3. Sea \(\varphi: \mathbb{N} \to \mathbb{N}\) la función definida inductivamente por \(\varphi(1) = 1\) y \(\varphi(n+1) = (n+1) \cdot \varphi(n)\). Entonces \(\varphi(1) = 1\), \(\varphi(2) = 2 \cdot 1\), \(\varphi(3) = 3 \cdot 2 \cdot 1\), etc. Así, \(\varphi(n) = n \cdot (n-1) \cdots 2 \cdot 1\), o sea, \(\varphi(n) = n!\) es el factorial de \(n\).

Ejemplo 4. Definamos \(f: \mathbb{N} \to \mathbb{Q}\) inductivamente poniendo
\[
f(1) = 1,\quad f(2) = 1/2 \text{ y, en general, } f(n+2) = \frac{f(n) + f(n+1)}{2}.
\]
Entonces \(f(3) = 3/4\), \(f(4) = 5/8\), etc. Cada valor \(f(n)\), para \(n > 2\), es la media aritmética de los dos valores anteriores de \(f\).

Ejemplo 5. La suma \(a_1 + \cdots + a_n\) de una \(n\)-upla de números naturales \(a_1, \ldots, a_n\) se define inductivamente así: \(a_1\) y \(a_1 + a_2\) ya se conocen. Suponiendo que ya se sabe sumar \(n\) números cualesquiera, se pone, por definición, \(a_1 + \cdots + a_{n+1} = (a_1 + \cdots + a_n) + a_{n+1}\). A primera vista, no se ve aquí la definición de una función \(f: \mathbb{N} \to X\), por inducción. En realidad, sin embargo, la suma de \(n\) números naturales (\(n\) fijo) es una función \(\sigma_n: \mathbb{N}^n \to \mathbb{N}\). Lo que se acaba de definir inductivamente fue la función \(f: n \to \sigma_n\). (¿Cuál es el codominio \(X\) de \(f\)?)

Ejemplo 6. También el producto \(a_1 \cdot a_2 \cdot \cdots \cdot a_n\) de una \(n\)-upla de números naturales se define inductivamente: \(a_1 \cdot a_2\) son conocidos. Se pone \(a_1 \cdot a_2 \cdot \cdots \cdot a_{n+1} = (a_1 \cdot a_2 \cdot \cdots \cdot a_n) \cdot a_{n+1}\).

\section{Conjuntos finitos e infinitos}

En este parágrafo, indicaremos con el símbolo \(I_n\) el conjunto \(\{1, \ldots, n\}\) de los números naturales desde 1 hasta \(n\). Más precisamente, dado \(n \in \mathbb{N}\), tenemos

\[
I_n = \{p \in \mathbb{N} \; ; \; 1 \leq p \leq n\}.
\]

Un conjunto \(X\) se llama finito cuando es vacío o cuando existe, para algún \(n \in \mathbb{N}\), una biyección
\[
\varphi: I_n \to X.
\]

En el primer caso, diremos que \(X\) tiene cero elementos. En el segundo caso, diremos que \(n \in \mathbb{N}\) es el número de elementos de \(X\), es decir, que \(X\) posee \(n\) elementos.

Los siguientes hechos se deducen inmediatamente de las definiciones:

a) cada conjunto \(I_n\) es finito y posee \(n\) elementos;\\
b) si \(f: X \to Y\) es una biyección, uno de esos conjuntos es finito si, y solo si, el otro lo es.

Intuitivamente, una biyección \(\varphi: I_n \to X\) significa un conteo de los elementos de \(X\). Poniendo \(\varphi(1) = x_1\), \(\varphi(2) = x_2, \dots, \varphi(n) = x_n\), tenemos \(X = \{x_1, x_2, \dots, x_n\}\). Esta es la representación ordinaria de un conjunto finito.

Para que el número de elementos de un conjunto no sea una noción ambigua, debemos probar que si existen dos biyecciones \(\varphi: I_n \to X\) y \(\psi: I_m \to X\), entonces \(m = n\). Considerando la función compuesta \(f = \psi^{-1} \circ \varphi: I_n \to I_m\), basta entonces probar que si existe una biyección \(f: I_n \to I_m\), entonces se tiene \(m = n\). Para fijar ideas, supongamos \(m \le n\). De ahí, \(I_m \subset I_n\). La unicidad del número de elementos de un conjunto finito será, por tanto, una consecuencia de la proposición más general siguiente:

\textbf{Teorema 3}. Sea \(A \subset I_n\). Si existe una biyección \(f: I_n \to A\), entonces \(A = I_n\).

\textbf{Demostración}. Usaremos inducción en \(n\). El resultado es obvio para \(n = 1\). Supongamos que sea válido para un cierto \(n\) y consideremos una biyección \(f: I_{n+1} \to A\). Pongamos \(a = f(n + 1)\). La restricción de \(f\) a \(I_n\) proporciona una biyección \(f': I_n \to A - \{a\}\). Si tenemos \(A - \{a\} \subset I_n\), entonces, por la hipótesis de inducción, concluiremos que \(A - \{a\} = I_n\), de donde \(a = n + 1\) y \(A = I_{n+1}\). Si, en cambio, no es \(A - \{a\} \subset I_n\), entonces debe tenerse \(n + 1 \in A - \{a\}\). En este caso, existe \(p \in I_{n+1}\) tal que \(f(p) = n + 1\). Entonces definiremos una nueva biyección \(g: I_{n+1} \to A\) poniendo \(g(x) = f(x)\) si \(x \neq p\) y \(x \neq n + 1\), mientras que \(g(p) = a\), \(g(n + 1) = n + 1\). Ahora, la restricción de \(g\) a \(I_n\) nos dará una biyección \(g'': I_n \to A - \{n + 1\}\). Evidentemente, \(A - \{n + 1\} \subset I_n\). Luego, por la hipótesis de inducción, \(A - \{n + 1\} = I_n\), de donde \(A = I_{n+1}\). Esto concluye la demostración.

\textbf{Corolario 1}. Si existe una biyección \(f: I_m \to I_n\) entonces \(m = n\). En consecuencia, si existen dos biyecciones \(\psi: I_n \to X\) y \(\varphi: I_m \to X\), debe tenerse \(m = n\).

En efecto, podemos suponer, para fijar las ideas, que \(m \le n\). Entonces \(I_m \subset I_n\). Tomando \(A = I_m\) en el Teorema 3, obtenemos \(I_m = I_n\) y, por tanto, \(m = n\).

\textbf{Corolario 2.} No puede existir una biyección \( f: X \to Y \) de un conjunto finito \( X \) sobre una parte propia \( Y \subset X \).

En efecto, siendo \( X \) finito, existe una biyección \( \varphi: I_n \to X \). Sea \( A = \varphi^{-1}(Y) \). Entonces \( A \) es una parte propia de \( I_n \) y la restricción de \( \varphi \) a \( A \) proporciona una biyección \( \varphi': A \to Y \).

\begin{center}
    \begin{tikzcd}
        X \arrow[r, "f"] & Y \\
        I_n \arrow[u, "\varphi"] \arrow[r, dashed, "g"] & \mathcal{A} \arrow[u, swap, "\varphi'"]
    \end{tikzcd}
\end{center}

La composición \( g = (\varphi')^{-1} \circ f \circ \varphi: I_n \to A \) sería entonces una biyección de \( I_n \) sobre su parte propia \( A \), lo que contradice el Teorema 3. Luego no existe la biyección \( f \).

\textbf{Teorema 4.} Si \( X \) es un conjunto finito entonces todo subconjunto \( Y \subset X \) es finito. El número de elementos de \( Y \) no excede al de \( X \) y solo es igual cuando \( Y = X \).

\textbf{Demostración.} Basta probar el teorema en el caso en que \( X = I_n \). Esto es claro para \( n = 1 \) pues las únicas partes de \( I_1 \) son \( \emptyset \) y \( I_1 \). Supongamos el teorema demostrado para \( I_n \) y consideremos un subconjunto \( Y \subset I_{n+1} \). Si \( Y \subset I_n \) entonces, por la hipótesis de inducción, \( Y \) será un conjunto finito cuyo número de elementos es \( \leq n \) y, por tanto, \( \leq n + 1 \). Si, por el contrario, \( n + 1 \in Y \) entonces \( Y - \{n + 1\} \subset I_n \) y consecuentemente (salvo en el caso trivial \( Y = \{n + 1\} \)) existe una biyección \( \psi: I_p \to Y - \{n + 1\} \), con \( p \leq n \). Definiremos entonces una biyección \( \varphi: I_{p+1} \to Y \), poniendo \( \varphi(x) = \psi(x) \) para \( x \in I_p \) y \( \varphi(p + 1) = n + 1 \). Se sigue que \( Y \) es finito y su número de elementos no excede \( p + 1 \). Como \( p \leq n \), tenemos \( p + 1 \leq n + 1 \). Resta solo mostrar que si \( Y \subset I_n \) tiene \( n \) elementos entonces \( Y = I_n \). Esto es claro pues, por el Corolario 2 del Teorema 3, no puede haber una biyección de \( I_n \) sobre una parte propia suya \( Y \).

\textbf{Corolario 1.} Sea \( f: X \to Y \) una función inyectiva. Si \( Y \) es finito entonces \( X \) también lo es. Además, el número de elementos de \( X \) no excede al de \( Y \).

En efecto, \( f \) define una biyección de \( X \) sobre su imagen \( f(X) \), la cual es finita por ser una parte del conjunto finito \( Y \). Además, el número de elementos de \( f(X) \), que es igual al de \( X \), no excede al de \( Y \).

\textbf{Corolario 2.} Sea \( g: X \to Y \) una función sobreyectiva. Si \( X \) es finito entonces \( Y \) también lo será y su número de elementos no excede al de \( X \).

En efecto, como vimos en el Capítulo I (§4), \( g \) posee una inversa por la derecha, es decir, existe una función \( f: Y \to X \) tal que \( g \circ f = \text{id}_Y \). Entonces \( g \) es inversa por la izquierda de \( f \) y, por tanto, \( f \) es una función inyectiva de \( Y \) en el conjunto finito \( X \). Se sigue del Corolario 1 que \( Y \) es finito y su número de elementos no excede al de \( X \).

Un conjunto \( X \) se llama infinito cuando no es finito. Más explícitamente, \( X \) es infinito cuando no es vacío y, además, cualquiera que sea \( n \in \mathbb{N} \), no existe una biyección \( \varphi: I_n \to X \).

Por ejemplo, el conjunto \( \mathbb{N} \) de los números naturales es infinito. En efecto, dada cualquier función \( \varphi: I_n \to \mathbb{N} \), con \( n > 1 \), sea \( p = \varphi(1) + \cdots + \varphi(n) \). Entonces \( p > \varphi(x) \) para todo \( x \in I_n \), de donde \( p \notin \varphi(I_n) \). Luego, ninguna función \( \varphi: I_n \to \mathbb{N} \) es sobreyectiva.

Otra manera de verificar que \( \mathbb{N} \) es infinito es considerar el conjunto \( P = \{2, 4, 6, \dots\} \) de los números pares y definir la biyección \( f: \mathbb{N} \to P \), donde \( f(n) = 2n \). Como \( P \) es una parte propia de \( \mathbb{N} \), se sigue del Corolario 2 del Teorema 3 que \( \mathbb{N} \) no es finito.

Los hechos que acabamos de establecer para conjuntos finitos proporcionan, por exclusión, resultados sobre conjuntos infinitos. Por ejemplo, si \( f: X \to Y \) es inyectiva y \( X \) es infinito, entonces \( Y \) también lo es. Dada \( f: X \to Y \) sobreyectiva, si \( Y \) es infinito, entonces \( X \) es infinito. O también, como acabamos de usar, si \( X \) admite una biyección sobre una de sus partes propias, entonces \( X \) es infinito.

Otros ejemplos de conjuntos infinitos son \( \mathbb{Z} \) y \( \mathbb{Q} \), pues ambos conjuntos contienes \( \mathbb{N} \). Según una proposición debida a Euclides, el conjunto de los números primos es infinito. Caracterizaremos ahora los subconjuntos finitos (y por tanto los infinitos) de \( \mathbb{N} \).

Un conjunto \( X \subset \mathbb{N} \) se llama acotado cuando existe un \( p \in \mathbb{N} \) tal que \( p \geq n \) para cualquiera que sea \( n \in X \).

\textbf{Teorema 5.} Sea \( X \subset \mathbb{N} \) no vacío. Las siguientes afirmaciones son equivalentes:

$(a)$ \( X \) es finito;\\
$(b)$ \( X \) es acotado;\\
$(c)$ \( X \) posee un mayor elemento.

\textbf{Demostración.} Probaremos que $(a) \Rightarrow (b)$, $(b) \Rightarrow (c)$ y $(c) \Rightarrow (a)$.

$(a) \Rightarrow (b)$ – Sea \( X = \{x_1, \dots, x_n\} \). Tomando \( p = x_1 + \dots + x_n \), tenemos \( p > x \), para todo \( x \in X \), luego \( X \) es acotado.

$(b) \Rightarrow (c)$ – Suponiendo \( X \subset \mathbb{N} \) acotado, se sigue que el conjunto \( A = \{ p \in \mathbb{N} ; p \geq n \text{ para todo } n \in X \} \) es no vacío. Por el Principio de Buena Ordenación, existe \( p_0 \in A \), que es el menor elemento de \( A \). Afirmamos que debe ser \( p_0 \in X \). En efecto, si fuese \( p_0 \notin X \) entonces tendríamos \( p_0 > n \) para todo \( n \in X \). Como \( X \neq \varnothing \), esto obligaría \( p_0 > 1 \), de donde \( p_0 = p_1 + 1 \). Si existiese algún \( n \in X \) con \( p_1 < n \), esto traería \( p_0 = p_1 + 1 \leq n \) y (como estamos suponiendo \( p_0 \notin X \)) \( p_0 < n \), un absurdo. Luego es \( p_1 \geq n \) para todo \( n \in X \). Pero esto significa \( p_1 \in A \), lo que es absurdo en vista de \( p_1 < p_0 \) y \( p_0 = \) menor elemento de \( A \). Por lo tanto debe ser \( p_0 \in X \). Como \( p_0 \geq n \) para todo \( n \in X \), concluimos que \( p_0 \) es el mayor elemento de \( X \).

$(c) \Rightarrow (a)$ – Si existe un elemento \( p \in X \) que es el mayor de todos, entonces \( X \) está contenido en \( I_p \) y, por consiguiente, \( X \) es finito, por el Teorema 4.

Un conjunto \( X \subset \mathbb{N} \) se llama ilimitado cuando no es acotado. Esto significa que, dado cualquier \( p \in \mathbb{N} \), existe algún \( n \in X \) tal que \( n > p \). Los conjuntos ilimitados \( X \subset \mathbb{N} \) son, como acabamos de ver, precisamente los subconjuntos infinitos de \( \mathbb{N} \).

\textbf{Teorema 6.} Sean \( X, Y \) conjuntos finitos disjuntos, con \( m \) y \( n \) elementos respectivamente. Entonces \( X \cup Y \) es finito y posee \( m + n \) elementos.

\textbf{Demostración.} Dadas las biyecciones \( \varphi: I_m \to X \) y \( \psi: I_n \to Y \), definamos la función \( \xi: I_{m+n} \to X \cup Y \) poniendo \( \xi(x) = \varphi(x) \) si \( 1 \leq x \leq m \) y \( \xi(m + x) = \psi(x) \) si \( 1 \leq x \leq n \). Como \( X \cap Y = \varnothing \), se comprueba inmediatamente que \( \xi \) es una biyección, lo que prueba el teorema.

\textbf{Corolario 1.} Sean \( X_1, \dots, X_k \) conjuntos finitos, dos a dos disjuntos, con \( m_1, \dots, m_k \) elementos respectivamente. Entonces \( X_1 \cup \dots \cup X_k \) es finito y posee \( m_1 + \dots + m_k \) elementos.

El corolario se obtiene aplicando el teorema \( k-1 \) veces.

\textbf{Corolario 2.} Sean \( Y_1, \dots, Y_k \) conjuntos finitos (no necesariamente disjuntos) con \( m_1, \dots, m_k \) elementos respectivamente. Entonces \( Y_1 \cup \dots \cup Y_k \) es finito y posee a lo sumo \( m_1 + \dots + m_k \) elementos.

Para cada \( i = 1, \dots, k \), sea \( X_i = \{(x, i) ; x \in Y_i\} \), es decir \( X_i = Y_i \times \{i\} \). Entonces los \( X_i \) son dos a dos disjuntos y cada \( X_i \) posee \( m_i \) elementos. Luego \( X_1 \cup \dots \cup X_k \) es finito y tiene \( m_1 + \dots + m_k \) elementos, por el Corolario 1. Ahora la aplicación \( f: X_1 \cup \dots \cup X_k \to Y_1 \cup \dots \cup Y_k \), definida por \( f(x, i) = x \), es sobreyectiva. El resultado se sigue.

\textbf{Corolario 3.} Sean \( X_1, \dots, X_k \) conjuntos finitos con \( m_1, \dots, m_k \) elementos respectivamente. El producto cartesiano \( X_1 \times \dots \times X_k \) es finito y posee \( m_1 \cdot m_2 \cdot \dots \cdot m_k \) elementos.

Basta demostrar el Corolario 3 para \( k = 2 \), pues el caso general se reduce a este mediante aplicaciones sucesivas del mismo resultado. Ahora, dados \( X \) e \( Y \) finitos, escribamos \( Y = \{y_1, \dots, y_n\} \). Entonces \( X \times Y = X_1 \cup \dots \cup X_n \), donde \( X_i = X \times \{y_i\} \). Como los \( X_i \) son dos a dos disjuntos y poseen el mismo número de elementos (digamos \( m \)) que \( X \), concluimos que \( X \times Y \) posee \( m + m + \dots + m = m \cdot n \) elementos.

\textbf{Corolario 4.} Si \( X \) e \( Y \) son conjuntos finitos y poseen respectivamente \( m \) y \( n \) elementos, entonces el conjunto \( F(X;Y) \) de todas las funciones \( f: X \to Y \) es finito y posee \( n^m \) elementos.

Consideremos primero el caso en que \( X = I_m \). Entonces una función \( f: I_m \to Y \) es simplemente una \( m \)-tupla de elementos de \( Y \). En otras palabras, \( F(I_m;Y) = Y \times \dots \times Y \) (\( m \) factores). Por el Corolario 3, el número de elementos de \( F(I_m;Y) \) es \( n^m \). El caso general se reduce a este. En efecto, sea \( \varphi: I_m \to X \) una biyección. La correspondencia que a cada \( f: X \to Y \) le asocia \( f \circ \varphi: I_m \to Y \) es una biyección de \( F(X;Y) \) sobre \( F(I_m;Y) \). Luego estos dos conjuntos poseen el mismo número de elementos, a saber, \( n^m \).

\section{Conjuntos enumerables}

Un conjunto \( X \) se dice enumerable cuando es finito o cuando existe una biyección \( f: \mathbb{N} \to X \). En el segundo caso, \( X \) se dice infinito enumerable y, poniendo \( x_1 = f(1), x_2 = f(2), \dots, x_n = f(n), \dots \), se tiene \( X = \{x_1, x_2, \dots, x_n, \dots\} \). Cada biyección \( f: \mathbb{N} \to X \) se llama una enumeración (de los elementos) de \( X \).

\textbf{Ejemplo 7.} La biyección \( f: \mathbb{N} \to P \), \( f(n) = 2^n \), muestra que el conjunto \( P \) de los números naturales pares es infinito enumerable. Análogamente, \( g: n \mapsto 2n-1 \) define una biyección de \( \mathbb{N} \) sobre el conjunto de los números naturales impares, el cual es, por tanto, infinito enumerable. También el conjunto \( \mathbb{Z} \) de los números enteros es enumerable. Basta notar que la función \( h: \mathbb{Z} \to \mathbb{N} \), definida por \( h(n) = 2n \) cuando \( n \) es positivo y \( h(n) = -2n + 1 \) cuando \( n \) es negativo o cero, es una biyección. Luego, \( h^{-1}: \mathbb{N} \to \mathbb{Z} \) es una enumeración de \( \mathbb{Z} \).

\textbf{Teorema 7.} Todo conjunto infinito \( X \) contiene un subconjunto infinito enumerable.

\textbf{Demostración.} Basta definir una función inyectiva \( f: \mathbb{N} \to X \). Para ello, comenzamos eligiendo, en cada subconjunto no vacío \( A \subset X \), un elemento \( x_A \in A \). A continuación, definimos \( f \) por inducción. Ponemos \( f(1) = x_X \) y, suponiendo ya definidos \( f(1), \dots, f(n) \), escribimos \( A_n = X - \{f(1), \dots, f(n)\} \). Como \( X \) no es finito, \( A_n \) no es vacío. Pondremos entonces \( f(n+1) = x_{A_n} \). Esto completa la definición inductiva de la función \( f: \mathbb{N} \to X \). Afirmamos que \( f \) es inyectiva. En efecto, dados \( m \neq n \) en \( \mathbb{N} \), sea \( m < n \). Entonces \( f(m) \in \{f(1), \dots, f(n-1)\} \) mientras que \( f(n) \in \complement \{f(1), \dots, f(n-1)\} \). Luego \( f(m) \neq f(n) \). La imagen \( f(\mathbb{N}) \) es, por tanto, un subconjunto infinito enumerable de \( X \).

\textbf{Corolario.} Un conjunto \( X \) es infinito si, y solo si, existe una biyección \( f: X \to Y \) de \( X \) sobre una parte propia \( Y \subset X \).

En efecto, si tal biyección existe, \( X \) será infinito, por el Corolario 2 del Teorema 3. Recíprocamente, si \( X \) es infinito, contiene un subconjunto infinito enumerable \( A = \{a_1, a_2, \dots, a_n, \dots\} \). Sea \( Y = (X - A) \cup \{a_2, a_4, \dots, a_{2n}, \dots\} \). Evidentemente, \( Y \) es una parte propia de \( X \). Definimos una biyección \( f: X \to Y \) poniendo \( f(x) = x \) si \( x \in X - A \) y \( f(a_n) = a_{2n} \).

Evidentemente, el corolario anterior también puede enunciarse así: "Un conjunto es finito si, y solo si, no admite una biyección sobre una parte propia suya". Se obtiene así una caracterización de los conjuntos finitos en la que no interviene el conjunto \( \mathbb{N} \). Esta fue la manera como Dedekind definió conjunto finito.

\textbf{Teorema 8.} Todo subconjunto \( X \subset \mathbb{N} \) es enumerable.

\textbf{Demostración.} Si \( X \) es finito, es enumerable. Si es infinito, definiremos inductivamente una biyección \( f: \mathbb{N} \to X \). Pondremos \( f(1) = \) menor elemento de \( X \). Supongamos \( f(1), \dots, f(n) \) definidos de modo que satisfagan las siguientes condiciones: (a) \( f(1) < f(2) < \dots < f(n) \); (b) poniendo \( B_n = X - \{f(1), \dots, f(n)\} \), se tiene \( f(n) < x \) para todo \( x \in B_n \). A continuación, notando que \( B_n \neq \varnothing \) (pues \( X \) es infinito), definimos \( f(n+1) = \) menor elemento de \( B_n \). Esto completa la definición de \( f: \mathbb{N} \to X \), de modo que se mantienen las condiciones (a) y (b) para todo \( n \in \mathbb{N} \). Se sigue de (a) que \( f \) es inyectiva. Por otro lado, (b) implica que \( f \) es sobreyectiva pues si existiese algún \( x \in X - f(\mathbb{N}) \), tendríamos \( x \in B_n \) para todo \( n \in \mathbb{N} \), por tanto, \( x > f(n) \), cualquiera que fuese \( n \in \mathbb{N} \). Entonces el conjunto infinito \( f(\mathbb{N}) \subset \mathbb{N} \) sería acotado, una contradicción, en vista del Teorema 5.

\textbf{Corolario 1.} Un subconjunto de un conjunto enumerable es enumerable. O sea: si \( f: X \to Y \) es inyectiva e \( Y \) es enumerable, entonces \( X \) es enumerable.

Una función \( f: X \to Y \), donde \( X, Y \subset \mathbb{N} \), se llama creciente cuando, dados \( m < n \) en \( X \), se tiene \( f(m) < f(n) \).

\textbf{Corolario 2.} Dado un subconjunto infinito \( X \subset \mathbb{N} \), existe una biyección creciente \( f: \mathbb{N} \to X \).

Esto es lo que se demostró en el Teorema 8.
Se sigue del Teorema 8 que el conjunto de los números primos es (infinito y) enumerable.

\textbf{Teorema 9.} Sea \( X \) un conjunto enumerable. Si \( f: X \to Y \) es sobreyectiva, entonces \( Y \) es enumerable.

\textbf{Demostración.} Existe \( g: Y \to X \) tal que \( f \circ g = id_Y \). Luego \( f \) es una inversa por la izquierda de \( g \), y, por tanto, \( g \) es inyectiva. Se sigue que \( Y \) es enumerable. (Corolario 1 del Teorema 8.)

\textbf{Teorema 10.} Sean \( X, Y \) conjuntos enumerables. El producto cartesiano \( X \times Y \) es enumerable.

\textbf{Demostración.} Existen funciones inyectivas \( \varphi: X \to \mathbb{N} \) y \( \psi: Y \to \mathbb{N} \). Luego \( g: X \times Y \to \mathbb{N} \times \mathbb{N} \), dada por \( g(x, y) = (\varphi(x), \psi(y)) \) es inyectiva. Siendo así, por el Corolario 1 del Teorema 8, basta probar que \( \mathbb{N} \times \mathbb{N} \) es enumerable. Para ello definimos la función \( f: \mathbb{N} \times \mathbb{N} \to \mathbb{N} \), donde \( f(m, n) = 2^m \cdot 3^n \). Por la unicidad de la descomposición en factores primos, \( f \) es inyectiva, de donde proporciona una biyección de \( \mathbb{N} \times \mathbb{N} \) sobre el conjunto enumerable \( f(\mathbb{N} \times \mathbb{N}) \subset \mathbb{N} \).

\textbf{Corolario 1.} El conjunto \( \mathbb{Q} \) de los números racionales es enumerable.

En efecto, si indicamos por \( \mathbb{Z}^* \) el conjunto de los números enteros \( \neq 0 \), veremos que \( \mathbb{Z}^* \) es enumerable. Luego es también enumerable el producto cartesiano \( \mathbb{Z} \times \mathbb{Z}^* \). Ahora, la función \( f: \mathbb{Z} \times \mathbb{Z}^* \to \mathbb{Q} \), definida por \( f(m, n) = \frac{m}{n} \), es sobreyectiva. Se sigue del Teorema 9 que \( \mathbb{Q} \) es enumerable.

\textbf{Corolario 2.} Sean \( X_1, X_2, \dots, X_n, \dots \) conjuntos enumerables.

La unión \( X = \bigcup_{n=1}^{\infty} X_n \) es enumerable.

En palabras: una unión enumerable de conjuntos enumerables es un conjunto enumerable.

Para demostrarlo, tenemos, para cada \( m \in \mathbb{N} \), una función sobreyectiva \( f_m: \mathbb{N} \to X_m \). A continuación, definamos una función \( f: \mathbb{N} \times \mathbb{N} \to X \) poniendo \( f(m, n) = f_m(n) \). Se ve inmediatamente que \( f \) es sobreyectiva. Como \( \mathbb{N} \times \mathbb{N} \) es enumerable, se concluye del Teorema 9 que \( X \) es enumerable.

En particular, una unión finita \( X = X_1 \cup \dots \cup X_n \) de conjuntos enumerables es enumerable: basta aplicar el corolario anterior, con \( X_{n+1} = X_{n+2} = \dots = \varnothing \).

Aplicaciones repetidas del Teorema 10 muestran que, si \( X_1, \dots, X_k \) son conjuntos enumerables, su producto cartesiano \( X = X_1 \times \dots \times X_k \) es enumerable. No es verdad, sin embargo, que el producto cartesiano \( X = \prod_{n=1}^{\infty} X_n \) de una secuencia de conjuntos enumerables sea siempre enumerable. Examinaremos este fenómeno en el párrafo siguiente.

\section{Conjuntos no enumerables}

El principal ejemplo de conjunto no enumerable que encontraremos en este libro será el conjunto \( \mathbb{R} \) de los números reales. Esto se probará en el capítulo siguiente. Aquí veremos, mediante un argumento simple debido a Cantor, que existen conjuntos no enumerables. Más generalmente, mostraremos que, dado cualquier conjunto \( X \), existe siempre un conjunto cuyo número cardinal es mayor que el de \( X \).

No definiremos qué es el número cardinal de un conjunto pero diremos que dos conjuntos \( X \) e \( Y \) tienen el mismo número cardinal, y escribiremos

\[ \text{card}(X) = \text{card}(Y), \]

para significar que existe una biyección \( f: X \to Y \).

Así, dos conjuntos finitos tienen el mismo número cardinal si, y solo si, poseen el mismo número de elementos de acuerdo con la definición dada en el §3. Si \( X \) es infinito enumerable, se tiene \( \text{card}(X) = \text{card}(Y) \) si, y solo si, \( Y \) es infinito enumerable.

Dados los conjuntos \( X, Y \), diremos que \( \text{card}(X) < \text{card}(Y) \) cuando exista una función inyectiva \( f: X \to Y \) pero no exista una función sobreyectiva \( f: X \to Y \).

El Teorema 7 muestra que, para todo conjunto infinito \( X \), se tiene \( \text{card}(\mathbb{N}) \leq \text{card}(X) \). Así, el número cardinal de un conjunto infinito enumerable es el menor de los números cardinales de los conjuntos infinitos.

Recordamos que, dados dos conjuntos \( X, Y \), el símbolo \( \mathcal{F}(X; Y) \) representa el conjunto de todas las funciones \( f: X \to Y \).

\textbf{Teorema 11.} (Cantor). Sean \( X \) un conjunto arbitrario e \( Y \) un conjunto que contenga al menos dos elementos. Ninguna función \( \varphi: X \to \mathcal{F}(X; Y) \) es sobreyectiva.

\textbf{Demostración.} Dada \( \varphi: X \to \mathcal{F}(X; Y) \), indicaremos con \( \varphi_x \) el valor de \( \varphi \) en el punto \( x \in X \). Así, \( \varphi_x \) es una función de \( X \) en \( Y \). Construiremos ahora una \( f \in \mathcal{F}(X; Y) \) tal que \( \varphi_x \neq f \) para todo \( x \in X \). Esto se hace eligiendo, para cada \( x \in X \), un elemento \( f(x) \in Y \), diferente de \( \varphi_x(x) \). Como \( Y \) contiene al menos dos elementos, esto es posible. La función \( f: X \to Y \) así obtenida es tal que \( f(x) \neq \varphi_x(x) \) y, por tanto, \( f \neq \varphi_x \), para todo \( x \in X \). Luego \( f \notin \varphi(X) \) y, en consecuencia, \( \varphi \) no es sobreyectiva.

\textbf{Corolario.} Sean \( X_1, X_2, \ldots, X_n, \ldots \) conjuntos infinitos enumerables. El producto cartesiano \( \prod_{n=1}^{\infty} X_n \) no es enumerable.

Basta considerar el caso en que todos los \( X_n \) son iguales a \( \mathbb{N} \). En este caso, \( \prod X_n = \mathcal{F}(\mathbb{N}; \mathbb{N}) \), que no es enumerable, por el Teorema 11.

El argumento usado en la demostración del Teorema 11 se llama "método de la diagonal, de Cantor". Este nombre se debe al caso particular en que \( X = \mathbb{N} \). Los elementos de \( \mathcal{F}(\mathbb{N}; Y) \) son sucesiones de elementos de \( Y \). Para probar que ninguna función \( \varphi : \mathbb{N} \to \mathcal{F}(\mathbb{N}; Y) \) es sobreyectiva, escribimos \( \varphi(1) = s_1 \), \( \varphi(2) = s_2 \), ... etc., donde \( s_1, s_2, \ldots \) son sucesiones de elementos de \( Y \). Luego

\[
s_1 = (y_{11}, y_{12}, y_{13}, \ldots)
\]
\[
s_2 = (y_{21}, y_{22}, y_{23}, \ldots)
\]
\[
s_3 = (y_{31}, y_{32}, y_{33}, \ldots)
\]
\begin{center}
    ......
\end{center}

A continuación, formamos una nueva sucesión \( s = (y_1, y_2, y_3, \ldots) \) de elementos de \( Y \) simplemente eligiendo, para cada \( n \in \mathbb{N} \), un elemento \( y_n \in Y \) diferente del \( n \)-ésimo término de la diagonal: \( y_n \neq y_{nn} \). La sucesión \( s \) no pertenece a la lista de las sucesiones \( s_n \) (pues el \( n \)-ésimo término de \( s \) es diferente del \( n \)-ésimo término de \( s_n \)). Así, ninguna lista enumerable puede agotar todas las funciones en \( \mathcal{F}(\mathbb{N}; Y) \).

Como en el caso particular de este teorema, consideremos el conjunto de todas las listas infinitas (enumerables) que se pueden formar utilizando solo los dígitos 0 y 1, como por ejemplo,

\[
0 \quad 0 \quad 0 \quad 0 \quad 0 \quad 0 \quad \cdots
\]
\[
1 \quad 1 \quad 0 \quad 0 \quad 0 \quad 1 \quad 1 \quad \cdots
\]
\[
0 \quad 1 \quad 0 \quad 1 \quad 1 \quad 1 \quad 0 \quad 1 \quad \cdots
\]

El conjunto de todas estas listas de ceros y unos no es enumerable. Sea \( \mathcal{P}(A) \) el conjunto de las partes de un conjunto dado \( A \). Considerando el conjunto de dos elementos \( \{0, 1\} \), veremos ahora que existe una biyección:

\[
\xi : \mathcal{P}(A) \to \mathcal{F}(A; \{0, 1\}).
\]

A cada \( X \in \mathcal{P}(A) \), es decir, a cada subconjunto \( X \subset A \), le asociamos la función \( \xi_X : A \to \{0, 1\} \) llamada función característica del conjunto \( X \): se tiene \( \xi_X(x) = 1 \) si \( x \in X \) y \( \xi_X(x) = 0 \) si \( x \notin X \).

La correspondencia \( X \mapsto \xi_X \) es una biyección de \( \mathcal{P}(A) \) sobre \( \mathcal{F}(A; \{0, 1\}) \). Su inversa asocia a cada función \( f : A \to \{0, 1\} \) el conjunto \( X \) de los puntos \( x \in A \) tales que \( f(x) = 1 \).

Como \( \{0, 1\} \) tiene dos elementos, se sigue del Teorema 11 que ninguna función \( \varphi : A \to \mathcal{F}(A; \{0, 1\}) \) es sobreyectiva. En consecuencia, ninguna función \( \psi : A \to \mathcal{P}(A) \) es sobreyectiva. (Si lo fuera, \( \varphi = \xi \circ \psi : A \to \mathcal{F}(A; \{0, 1\}) \) también sería sobreyectiva.)

Pero existe una función inyectiva evidente \( f : A \to \mathcal{P}(A) \), definida por \( f(x) = \{x\} \). Concluimos entonces que \( \text{card}(A) < \text{card}[\mathcal{P}(A)] \), para todo conjunto \( A \).

Sobre números cardinales de conjuntos, informamos que, dados dos conjuntos \( A \) y \( B \) cualesquiera, vale una, y solo una, de las siguientes alternativas: \( \text{card}(A) = \text{card}(B) \), \( \text{card}(A) < \text{card}(B) \) o \( \text{card}(B) < \text{card}(A) \). Además, si existen una función inyectiva \( f : A \to B \) y una función inyectiva \( g : B \to A \), existirá también una biyección \( h : A \to B \).

\textbf{Ejercicios}

1. Pruebe que, en presencia de los axiomas P1 y P2, el axioma (A) siguiente es equivalente a P3. (A) Para todo subconjunto no vacío \( A \subset \mathbb{N} \), se tiene \( A - s(A) \neq \varnothing \).

2. Dados los números naturales \( a \), \( b \), pruebe que existe un número natural \( m \) tal que \( m \cdot a > b \).

3. Sea \( a \) un número natural. Si un conjunto \( X \) es tal que \( a \in X \) y, además, \( n \in X \Rightarrow n + 1 \in X \), entonces \( X \) contiene todos los números naturales \( \geq a \).

4. Intente descubrir, independientemente, algunas de las demostraciones omitidas en el texto. Si no logra alguna, consulte uno de los libros aquí citados, u otros de su predilección. [Sugerencia: Prácticamente todas las proposiciones sobre \( \mathbb{N} \) se demuestran por inducción.]

5. Un elemento \( a \in \mathbb{N} \) se llama antecesor de \( b \in \mathbb{N} \) cuando se tiene \( a < b \) pero no existe \( c \in \mathbb{N} \) tal que \( a < c < b \). Pruebe que, excepto 1, todo número natural posee un antecesor.

6. Use inducción para demostrar los siguientes hechos:

\begin{itemize}
    \item[a)] \( 2(1 + 2 + \cdots + n) = n(n+1) \)
    \item[b)] \( 1 + 3 + 5 + \cdots + (2n+1) = (n+1)^2 \)
    \item[c)] \( (a-1)(1 + a + \cdots + a^n) = a^{n+1} - 1 \), cualesquiera sean \( a, n \in \mathbb{N} \)
    \item[d)] \( n \geq 4 \Rightarrow n! > 2^n \)
\end{itemize}

7. Use el Segundo Principio de Inducción para demostrar la unicidad de la descomposición de un número natural en factores primos.

8. Sea \( X \) un conjunto con \( n \) elementos. Use inducción para probar que el conjunto de las biyecciones (o permutaciones) \( f: X \to X \) tiene \( n! \) elementos.

9. Sean \( X \) e \( Y \) conjuntos finitos.

\begin{itemize}
    \item[a)] Pruebe que \( \text{card}(X \cup Y) + \text{card}(X \cap Y) = \text{card}(X) + \text{card}(Y) \).
    \item[b)] ¿Cuál sería la fórmula correspondiente para tres conjuntos?
    \item[c)] Generalice.
\end{itemize}

10. Dado un conjunto finito \( X \), pruebe que una función \( f: X \to X \) es inyectiva si, y solo si, es sobreyectiva (y por tanto una biyección).

11. Formula matemáticamente y demuestra el siguiente hecho (conocido como el "principio de los cajones" o principio de palomar). Si \(m < n\), entonces, de cualquier modo que se guarden \(n\) objetos en \(m\) cajones, habrá siempre al menos un cajón que contendrá más de un objeto.

12. Sea \(X\) un conjunto con \(n\) elementos. Determina el número de funciones inyectivas \(f: I_p \to X\).

13. ¿Cuántos subconjuntos con \(p\) elementos posee un conjunto \(X\), sabiendo que \(X\) tiene \(n\) elementos?

14. Demuestra que si \(A\) tiene \(n\) elementos, entonces \(\mathcal{P}(A)\) tiene \(2^n\) elementos.

15. Define una función sobreyectiva \(f: \mathbb{N} \to \mathbb{N}\) tal que, para todo \(n \in \mathbb{N}\), el conjunto \(f^{-1}(n)\) sea infinito.

16. Demuestra que si \(X\) es infinito enumerable, el conjunto de las partes finitas de \(X\) también es (infinito) enumerable.

17. Sea \(f: X \to X\) una función. Un subconjunto \(Y \subset X\) se llama estable respecto a \(f\) cuando \(f(Y) \subset Y\). Demuestra que un conjunto \(X\) es finito si, y solo si, existe una función \(f: X \to X\) que solo admite los subconjuntos estables \(\varnothing\) y \(X\).

18. Sea \(f: X \to X\) una función inyectiva tal que \(f(X) \neq X\). Tomando \(x \in X - f(X)\), demuestra que los elementos \(x\), \(f(x)\), \(f(f(x))\), ... son dos a dos distintos.

19. Sean \(X\) un conjunto infinito e \(Y\) un conjunto finito. Muestra que existe una función sobreyectiva \(f: X \to Y\) y una función inyectiva \(g: Y \to X\).

20.a) Si \(X\) es finito e \(Y\) es enumerable, entonces \(\mathcal{F}(X; Y)\) es enumerable.\\
b) Para cada función \(f: \mathbb{N} \to \mathbb{N}\) sea \(A_f = \{ n \in \mathbb{N} ; f(n) \neq 1 \}\). Demuestra que el conjunto \(X\) de las funciones \(f: \mathbb{N} \to \mathbb{N}\) tales que \(A_f\) es finito es un conjunto enumerable.

21. Obtén una descomposición \( \mathbb{N} = X_1 \cup X_2 \cup \cdots \cup X_n \cup \ldots \) tal que los conjuntos \( X_1, X_2, \ldots, X_n, \ldots \) sean infinitos y dos a dos disjuntos.

22. Define \( f: \mathbb{N} \times \mathbb{N} \to \mathbb{N} \), poniendo \( f(1,n) = 2n - 1 \) y \( f(m+1,n) = 2^m \cdot (2n - 1) \). Prueba que \( f \) es una biyección.

23. Sea \( X \subset \mathbb{N} \) un subconjunto infinito. Demuestra que existe una única biyección creciente \( f: \mathbb{N} \to X \).

24. Demuestra que todo conjunto infinito se descompone como unión de una infinidad enumerable de conjuntos infinitos, dos a dos disjuntos.

25. Sea \( A \) un conjunto. Dadas dos funciones \( f, g: A \to \mathbb{N} \), define la suma \( f+g: A \to \mathbb{N} \), el producto \( f \cdot g: A \to \mathbb{N} \), y da el significado de la afirmación \( f \leq g \). Indicando con \( \xi_X \) la función característica de un subconjunto \( X \subset A \), prueba:
\begin{itemize}
    \item[a)] \( \xi_{X \cap Y} = \xi_X \cdot \xi_Y \);  
    \item[b)] \( \xi_{X \cup Y} = \xi_X + \xi_Y - \xi_{X \cap Y} \). En particular, \( \xi_{X \cup Y} = \xi_X + \xi_Y \Leftrightarrow X \cap Y = \varnothing \);
    \item[c)] \( X \subset Y \Leftrightarrow \xi_X \leq \xi_Y \);
    \item[d)] \( \xi_{A-X} = 1 - \xi_X \).
\end{itemize}

26. Demuestra que el conjunto de las sucesiones crecientes \( (n_1 < n_2 < n_3 < \ldots) \) de números naturales no es enumerable.

27. Sean \( (\mathbb{N}, s) \) y \( (\mathbb{N}', s') \) dos pares formados, cada uno, por un conjunto y una función. Supongamos que ambos cumplen los axiomas de Peano. Demuestra que existe una única biyección \( f: \mathbb{N} \to \mathbb{N}' \) tal que \( f(1) = 1' \), \( f(s(n)) = s'(f(n)) \). Concluye que:
\begin{itemize}
    \item[a)] \( m < n \Leftrightarrow f(m) < f(n) \);
    \item[b)] \( f(m+n) = f(m) + f(n) \); y
    \item[c)] \( f(m \cdot n) = f(m) \cdot f(n) \).
\end{itemize}

28. Dada una sucesión de conjuntos \( A_1, A_2, \dots, A_n, \dots \), considere los conjuntos  
\[
\limsup A_n = \bigcap_{n=1}^{\infty} \left( \bigcup_{i=n}^{\infty} A_i \right) \quad \text{y} \quad \liminf A_n = \bigcup_{n=1}^{\infty} \left( \bigcap_{i=n}^{\infty} A_i \right).
\]
\begin{itemize}
    \item[a)] Demuestre que \(\limsup A_n\) es el conjunto de los elementos que pertenecen a \(A_n\) para una infinidad de valores de \(n\) y que \(\liminf A_n\) es el conjunto de los elementos que pertenecen a todo \(A_n\) salvo para un número finito de valores de \(n\).
    \item[b)] Concluya que \(\liminf A_n \subset \limsup A_n\).
    \item[c)] Muestre que si \(A_n \subset A_{n+1}\) para todo \(n\), entonces \(\liminf A_n = \limsup A_n = \bigcup_{n=1}^{\infty} A_n\).
    \item[d)] Por otro lado, si \(A_n \supset A_{n+1}\) para todo \(n\), entonces \(\liminf A_n = \limsup A_n = \bigcap_{n=1}^{\infty} A_n\).
    \item[e)] Dé un ejemplo de una sucesión \((A_n)\) tal que \(\limsup A_n \neq \liminf A_n\).
    \item[f)] Dé un ejemplo de una sucesión para la cual los dos límites coinciden pero \(A_m \not\subset A_n\) cualesquiera que sean \(m\) y \(n\).
\end{itemize}

29. Dados los conjuntos \(A\) y \(B\), suponga que existan funciones inyectivas \(f: A \to B\) y \(g: B \to A\). Demuestre que existe una biyección \(h: A \to B\). (Teorema de Cantor-Bernstein-Schröder.)
